% ==============================================================================
% DOKUMENTKODIERUNG, SPRACHE & SCHRIFT
% ==============================================================================
\usepackage[utf8]{inputenc}             % Ermöglicht UTF-8-Kodierung im Quelltext (bei pdfLaTeX nötig)
\usepackage[T1]{fontenc}                % Verwendet T1-Schriftkodierung für korrekte Trennung/Umlaute
\usepackage{lmodern}                    % Latin-Modern-Schriften passend zu T1
\usepackage[ngerman]{babel}             % Deutsche Sprache: Silbentrennung, Bezeichnungen, Typografie
\usepackage{microtype}                  % Allein das Einbinden des Paketes verbessert den Blocksatz optich, also ein must have
\usepackage{csquotes}                   % Sprachabhängige Anführungszeichen
\usepackage{textcase}                   % Kontrollierte Groß-/Kleinschreibung in beweglichen Argumenten

% ==============================================================================
% SEITENLAYOUT & STRUKTUR
% ==============================================================================
\usepackage{scrlayer-scrpage}           % Kopf- und Fußzeilen (KOMA-Script)
\usepackage[final]{pdfpages}            % Einbinden externer PDF-Seiten | Wichtig für die Messdaten
%\usepackage{subfiles}                   % Essentiel beim aufteilen in viele .tex, ermöglich auch da saubere referenzierung und einzelne kompilierung
%\usepackage{import}                     % Fügt neuen befehl hinzu, dass .tex dateien im unterordner nicht die main.tex als referenzpunkt für pfade nehmen
% Falls ich es umändere, dass jede Sektion eine .tex Datei ist, sind diese beiden Pakete definitiv von relevanz

% ==============================================================================
% FARBEN & BOXEN
% ==============================================================================
\usepackage[x11names]{xcolor}           % Erweiterte Farbdefinitionen (X11-Farbnamen)
\usepackage{tcolorbox}                  % Farbige Boxen für Text, Theoreme etc.
    \tcbuselibrary{most}                % Lädt nahezu alle tcolorbox-Erweiterungen
\usepackage{framed}                     % Einfache Rahmen um Text

% ==============================================================================
% TABELLEN & LISTEN
% ==============================================================================
\usepackage{tabularray}                 % Eine Sammlung an Tabellenverbessurngspaketen
    \UseTblrLibrary{booktabs}               % Ermöglicht die booktabs Tabellen zu haben
\usepackage{multicol}                   % Zellen über mehrere Spalten
\usepackage{multirow}                   % Zellen über mehrere Zeilen
\usepackage{enumitem}                   % Kontrolle über Listen (Abstände, Labels)

% ==============================================================================
% ABBILDUNGEN & GLEITOBJEKTE
% ==============================================================================
\usepackage{float}                      % Erzwingt Platzierung von Gleitobjekten mit [H]
\usepackage[section]{placeins}          % Ermöglicht das Platzieren von float barriers
\usepackage{graphicx}                   % Ermöglicht das Einbinden von Grafiken in PDF, PNG usw. Form
\usepackage{caption}                    % Anpassung von Beschriftungen (Abbildungen/Tabellen)
%Reihenfolge Einhalten!
\usepackage{subcaption}                 % Teilabbildungen mit eigenen Beschriftungen | Modernere Version von subfigures

% ==============================================================================
% MATHEMATIK, CHEMIE, PHYSIK & EINHEITEN
% ==============================================================================
\usepackage{amsmath, amsthm, amssymb}   % Mathematische Umgebungen, Symbole, Theoreme
\usepackage{mathtools}                  % Erweiterung für amsmath
\usepackage{derivative}                 % Fügt starke Makros für Ableitungen hinzu | Nicht nötig, falls physics aktiv ist
\usepackage{bm}                         % Bessere Notation für Vektoren
\usepackage{siunitx}                    % Einheitliche Darstellung von Zahlen und SI-Einheiten | sehr wichtig für das AP
\usepackage{upgreek}                    % Aufrechte griechische Buchstaben im Mathematikmodus
\usepackage{cancel}                     % Durchstreichen von mathematischen Ausdrücken
\usepackage{empheq}                     % Hervorheben und Einrahmen von Gleichungssystemen
\usepackage{chemformula}                % Unterstützung für chemischen Satz | Nur selten benötigt, entkommentieren falls benötigt.

%Physics ist nicht so gut, da es viele Konflikte mit Paketen hat!
%\usepackage[arrowdel]{physics}          % Kurzbefehle für physikalische Notation (Ableitungen, Bra-Ket, ...)
%\let\qty\relax                          %\qty von physics entfernen, da siunitx es auch verwendet

% ==============================================================================
% GRAFIKEN & PLOTS
% ==============================================================================
\usepackage{tikz}                       % Zeichnen von Vektorgrafiken
    \usepackage{scalerel}                   % Relatives Skalieren von Symbolen
    \usepackage{pict2e}                     % Erweiterte picture-Umgebung
    \usepackage{tkz-euclide}                % Geometrische Konstruktionen (baut auf TikZ auf)
    \usetikzlibrary{calc}                   % Koordinatenberechnungen
    \usetikzlibrary{patterns,arrows.meta}   % Musterfüllungen und moderne Pfeile
    \usetikzlibrary{shadows}                % Schattierungen
    \usetikzlibrary{external}               % Externes Kompilieren von TikZ-Grafiken

\usepackage{pgfplots}                   % Funktions- und Datenplots über TikZ
    \pgfplotsset{compat=newest}             % Neueste PGFPlots-Funktionalität
    \usepgfplotslibrary{statistics}         % Statistikplots
    \usepgfplotslibrary{fillbetween}        % Flächen zwischen Kurven
    \usepgfplotslibrary{colorbrewer}        % Ein Farbpaket für Kurven und Messpunkte, sodass keine eigenen Farben ausgewählt werden müssen

% ==============================================================================
% PROGRAMMIERUNG & EIGENE MAKROS
% ==============================================================================
\usepackage{listings}                   % Darstellung von Quellcode
\usepackage{environ}                    % Definieren eigener Umgebungen mit Zugriff auf deren Inhalt
\usepackage{xparse}                     % Verbessert das Erstellen von eigenen Macros

% ==============================================================================
% LITERATUR, LINKS & QUERVERWEISE
% ==============================================================================
\usepackage[
    backend=biber,                      % Standart backend fürs Zitieren | Pflicht, sollte immer als Option drin sein 
    style=numeric-comp,                 % Kompakte Numerische Zitierung, anstatt Author + Jahr | \cite{A} -> [1]; \cite{A,B,C} -> [1-3]
    sorting=none                        % Nicht alphabetisch Sortieren. Sondern das erste Zitat im Dokument ist auch die Nr. 1 usw.
    ]
    {biblatex}                          % Paket zur erzeugung von Literaturverzeichnissen | Einfach Quellen in literatur.bib machen
\usepackage{xurl}                       % Sorgt für schönere url

% Hyperref immer vor cleveref!
\usepackage{hyperref}                   % Hyperlinks und PDF-Metadaten (immer zuletzt laden)
%Cleveref als letztes
\usepackage[noabbrev, nameinlink, ngerman]{cleveref}         % Intelligente Querverweise (Abb., Gl., Tab. automatisch)
